% Section 3.7, from the summary table of lectures/L03/appendix/a_theory.md and from
% lectures/L03/README.md: the objectives, the evaluation questions and the next lesson.

\section{Sammanfattning}
\label{c3:sec:sammanfattning}

\begin{booktable}{@{}p{10em}L@{}}
  \thead{Begrepp} & \thead{Nyckelregel} \\
  \midrule
  Linjär ekvation & Isolera den obekanta med additions- och multiplikationsprincipen \\
  Proportionsekvation & Korsvis multiplikation: $\frac{a}{b} = \frac{c}{d} \Rightarrow ad = bc$ \\
  Olikhet & Samma regler som för en ekvation, men tecknet vänder vid multiplikation eller
    division med ett negativt tal \\
  Absolutbelopp & $\abs{ax + b} = c \Rightarrow$ två fall: $ax + b = c$ eller $ax + b = -c$ \\
\end{booktable}

\needspace{6\baselineskip}
\subsection*{Det här ska du kunna}
Efter det här kapitlet ska du:
\begin{itemize}
  \item Kunna lösa linjära ekvationer algebraiskt.
  \item Förstå och lösa olikheter.
  \item Förstå och beräkna absolutbeloppet av ett tal eller uttryck.
  \item Kunna ställa upp och lösa ekvationer och olikheter utifrån elektriska samband.
\end{itemize}

\testadigsjalv
\begin{enumerate}
  \item Lös ekvationen $4(x - 2) = 2x + 8$ och kontrollera svaret.
  \item Spänningen är $U = 3{,}6\,\text{V}$ och motståndet $R = 120\,\Omega$. Beräkna strömmen $I$
    med Ohms lag.
  \item Klämspänningen för ett belastat batteri är $U = 9 - 0{,}5I$. Bestäm de strömmar som
    uppfyller $U \geq 6\,\text{V}$.
\end{enumerate}

\needspace{6\baselineskip}
\subsection*{Nästa kapitel}
\Kapref{4} handlar om ekvationssystem: två ekvationer med två obekanta, som löses med
substitutionsmetoden eller additionsmetoden.
